\subsection{凸函数的逼近}

命题 5. --- 设 X 是局部凸空间 E 中的一个闭凸集。则每个在 X 上定义的下半连续凸函数 f 都是 E 中连续仿射线性函数限制到 X 上所得函数组成的一个族的上包络。

因为，集 \(A \subset E \times R\) （它由那些点 \((x, t)\) 组成：\(x \in X\) 且 \(t \geqslant f(x)\) ）是凸的 (II, 第 17 页, prop. 19) 且闭的，因为函数 \((x, t) \mapsto f(x) - t\) 是下半连续的。设 x 是 X 的任一点，且设 \(a \in R\) 满足 \(a < f(x)\) 。由 II, 第 38 页的 cor. 1，存在 \(E \times R\) 中的一个闭超平面 H，它包含 \((x, a)\) 且不与 A 相交。由于 \(E \times R\) 上每个连续线性型都形如

\[
(z, t) \rightarrow u (z) + \lambda t,
\]

其中 \(\lambda\in R\) 而 u 是 E 上的一个连续线性型，由此可知 H 有形如 \(u(z)+\lambda t=\alpha\) 的方程，且由于 H 包含 \((x,a)\) ，我们有 \(\alpha=u(x)+\lambda a\) 。而点 \((x,f(x))\in\mathbf{A}\) 不属于 H，因此 \(\lambda\neq0\) 。必要时除以 \(-\lambda\) ，可把 H 的方程写成 \(t-a=u(z-x)\) 。由于 \(f(x)-a>0\) ，我们有 \(f(z)>u(z-x)+a\) （对所有 \(z\in X\) ），这就证明了命题。

注记. --- 1) 由 prop. 5 推出，f 是一个递增有向函数族的上包络，该族的函数是 E 中连续且凸的函数限制到 X 上的限制。

\begin{enumerate}
\def\labelenumi{\arabic{enumi})}
\setcounter{enumi}{1}
\tightlist
\item
  进一步假设 X 是一个以 0 为顶点的闭凸锥，且 f 是正齐次的。则 f 是 E 中连续线性型限制到 X 上所得函数组成的一个族的上包络。因为，设 \((u_{\alpha})\) 是 E 中连续仿射线性函数组成的一个族，它们在 X 上的限制以 f 为上包络。记 \(u_{\alpha} = v_{\alpha} + \lambda_{\alpha}\) ，其中 \(\lambda_{\alpha} \in R\) ，而 \(v_{\alpha}\) 是 E 中的一个连续线性型。我们有 \(\lambda_{\alpha} = u_{\alpha}(0) \leqslant f(0) = 0\) 。另一方面，若 \(x \in X\) ，则对每个 \(\mu > 0\) 有
\end{enumerate}

\[
\mu^ {- 1} \lambda_ {\alpha} + v _ {\alpha} (x) = \mu^ {- 1} (\lambda_ {\alpha} + v _ {\alpha} (\mu x)) = \mu^ {- 1} u _ {\alpha} (\mu x) \leqslant \mu^ {- 1} f (\mu x) = f (x)
\]

因此在 X 中 \(u_{\alpha} \leqslant v_{\alpha} \leqslant f\) ，从而 \(f\) 是诸 \(v_{\alpha}\) 的上包络。

\begin{enumerate}
\def\labelenumi{\arabic{enumi})}
\setcounter{enumi}{2}
\tightlist
\item
  E 中连续仿射函数在 X 上的限制是 X 上的一个仿射函数（即既凹又凸，II, 第 17 页）；但可能发生这样的情形：在紧凸集 \(\mathbf{X} \subset \mathbf{E}\) 中存在连续仿射函数，它们不是 E 中连续仿射函数在 X 上的限制 (II, 第 78 页, exerc. II, c))。然而：
\end{enumerate}

命题 6. --- 设 f 是 Hausdorff 局部凸空间 E 的紧凸集 X 上的一个上半连续仿射函数。设 L 是 E 中连续仿射函数在 X 上的限制所成的集；则由那些 \(h \in L\) （它们满足 \(h(x) > f(x)\) ，对所有 \(x \in X\) ）组成的集 L' 是递减有向的，且其下包络等于 f。

我们可以假设 X 非空。设 u, v 是 L 的两个元，使得 \(u(x) > f(x)\) 且 \(v(x) > f(x)\) （对所有 \(x \in X\) ），并设 b 是 u 和 v 的一个上界常数。设 U（相应地 V）是由点 \((x, t)\) （在 \(X \times R\) 中满足 \(u(x) \leqslant t \leqslant b\) （相应地 \(v(x) \leqslant t \leqslant b\) ））所成的紧凸集，并设 F 是由 \((x, t) \in X \times R\) （满足 \(t \leqslant f(x)\) ）所成的集；F 在 \(X \times R\) 中是凸的且闭的。\(U \cup V\) 的凸包 K 不与 F 相交，因为 \(U \cup V\) 含于由 \((x, t) \in X \times R\) （满足 \(f(x) < t\) ）所成的集中，而后者是凸的且不与 F 相交。由于 K 是紧的 (II, 第 14 页, prop. 15)，我们可以用 \(E \times R\) 中的一个闭超平面 H 把 F 与 K 严格分离。对每个 \(x \in X\) ，超平面 H 把 \((x, f(x))\) 与 \((x, b)\) 严格分离，因此与直线 \(\{x\} \times R\) 交于单独一点 \(w(x)\) ；于是 H 是一个连续仿射函数的图象，其在 X 上的限制 w 属于 L，是 u 和 v 的一个下界，且满足不等式 \(w(x) > f(x)\) （对所有 \(x \in X\) ）。这就证明集 \(L'\) 是递减有向的。把 II, 第 39 页的 prop. 5 应用于 -f，可知 f 是 \(L'\) 的下包络。

推论. --- 设 \(f\) 是 \(\mathbf{X}\) 中的一个连续仿射函数；则存在一个序列 \((h_n)\) （由 \(\mathbf{L}\) 的元素组成），它一致收敛于 \(f\) （在 \(\mathbf{X}\) 中） 。

因为，prop. 6 与 Dini 定理 (GT, X, § 4.1, th. 1) 表明，对所有 \(n\) 存在 \(h_n \in \mathbf{L}\) 使 \(f \leqslant h_n \leqslant f + 1/n\) 。

